\section{Background and interpretation}\label{sec:context}
The quantitative issue for ordered matrices is the dependence on the edit
parameter.  The row and column orders are independent: an occurrence of a
pattern $F\in\{0,1\}^{s\times t}$ is a pair of increasing injections
$[s]\hookrightarrow[n]$ and $[t]\hookrightarrow[n]$, and all $st$ entries,
including the zeros, are prescribed. The corresponding hereditary
property consists of matrices avoiding $F$ as an induced ordered
submatrix.  A removal bound supplies a lower bound on the number of
occurrences whenever one must change a prescribed fraction of all
entries to obtain a member of this hereditary property.

For an individual fixed $F$, a polynomial bound means the existence of
$c_F,C_F>0$ such that the right side of \eqref{eq:removal-bound} holds
for \emph{every} size and \emph{every} matrix.  Failure is witnessed by
one fixed pattern and a sequence of arbitrarily large hosts.  The
quantifiers matter: choosing a different forbidden pattern as the
accuracy parameter decreases would address a different question.
The present family uses the same $H_4$ at every depth; the pattern is
independent of $h$, $m$ and $n_h$.


We normalize occurrences by $n^{2k}$, as in
\eqref{eq:removal-bound}, and count increasing injections, so a copy
has no additional factorial multiplicity. For a fixed $k$, using
$\binom{n}{k}^2$ instead changes the density only by a factor bounded
above and below by positive constants for $n\ge k$. Thus the failure
of a polynomial bound is independent of these two common conventions.

Alon, Ben-Eliezer and Fischer \cite{ABF} proved qualitative testing
results for hereditary properties of ordered graphs and matrices.
Alon and Ben-Eliezer \cite{ABE20} obtained a polynomial dependence for
a \emph{hitting} formulation: a small set of entries may intersect all
copies present in the original matrix. Hitting the original copies
need not produce a free matrix after entry changes, because a changed
zero or one can complete a new induced copy.  The relevant full removal
problem minimizes changes over the resulting, genuinely $F$-free
matrices.  This distinction is particularly significant for an induced
pattern, whose constraints include prescribed zeros as well as ones.
Alon and Ben-Eliezer also establish a full polynomial removal lemma
for every fixed finite family of binary matrices closed under row
permutations \cite[Theorem~1.5]{ABE20}, and hence, by transposition,
for families closed under column permutations. The singleton
$\{H_4\}$ has neither closure property, so this positive result does
not apply to the present obstruction. Gishboliner and Shapira
\cite[Conjecture~4.5]{GS25} formulate the full polynomial removal
conjecture for a single fixed ordered binary matrix.
Towsner\'s ordered-hypergraph removal lemma provides a further
point of comparison at the level of qualitative removal \cite{Towsner}.
Earlier work on unordered bipartite matrices and conditional regularity
provided the polynomial-query comparison \cite{AFN07}; the ordering
requirement removes the row-and-column permutation freedom used in that
setting.

The 2026 construction of \cite{OAI} gave a fixed $66\times66$ binary
obstruction. Its tree-based construction showed that forcing gadgets can
make edit distance persist over successive levels, while an order
constraint can confine complete copies to terminal locations.  Our
pattern \eqref{eq:pattern} has four rows and four columns; its top-left
$2\times2$ anti-diagonal marker and the two identity coupling blocks
provide a smaller frame for the same kind of propagation.  The
sharper size is accompanied here by exact distance and copy-count
identities.  We isolate the roles of the order, the two signs, and
terminal identification in subsequent sections.

The tree-based propagation and terminal-sign identification build on
\cite{OAI}. The additional results here are the explicit $4\times4$
pattern, its exact physical copy count, its optimal unrestricted
repair and the formal certification of these identities.

\subsection{Three quantities that must be kept separate}
Let $\mathcal C_H(A)$ be the set of increasing row/column tuples
realizing $H$. Define the \emph{transversal number}
\[
 \tau_H(A)=\min\{|S|: S\subseteq[n]\times[n],\;
                 S\cap\operatorname{cells}(Q)\ne\varnothing
                 \ \forall Q\in\mathcal C_H(A)\}.
\]
Every edit set turning $A$ into an $H$-free matrix intersects every
copy originally in $A$; consequently
\begin{equation}\label{eq:hitting_le}
 \tau_H(A)\le D_H(A).
\end{equation}
Equality does not follow from the definition: a hitting set has no
obligation to avoid copies created by its own edits.  For our host,
the protected-selection argument establishes the stronger lower bound
$D_{H_4}(A_h)\ge m^2$ \emph{for all resulting matrices}.  The
canonical repair gives a matching upper bound by an explicit final
$H_4$-free matrix.  Neither step is supplied merely by counting
original occurrences.

Indeed, the terminal classification below shows a concrete gap between
the two quantities. The $m$ leaf-diagonal cells
$\{(R_h(p),C_h(p)):0\le p<m\}$ intersect every original copy,
so $\tau_{H_4}(A_h)\le m<m^2=D_{H_4}(A_h)$.
Corollary~\ref{cor:exacttransversal} proves that the hitting bound is exact.
Every terminal copy contains its same-leaf diagonal as its
third-row, third-column cell. This is the hitting-versus-removal
phenomenon emphasized in \cite{OAI}, here retained by the smaller
pattern; it is not a new general separation principle.

It is also useful to distinguish physical copies from type-level
copies. In a block-constant host, quotient rows and columns encode
signatures and ordering, whereas block multiplicities control the
physical number of realizations. The quotient alone has no absolute
copy count unless its tuples are weighted by those multiplicities.
This separation is essential in the terminal enumeration.

\subsection{The four-by-four pattern as a forcing frame}
Partition $H_4$ after the second row and the second column. The
four component blocks are respectively $J_2-I_2$, $I_2$, $I_2$,
and the asymmetric body $P$. Thus two marker rows and columns
fix twelve entries of a potential induced copy; the residual body
has only four entries. If the residual body has the form
$\left(\begin{smallmatrix}a&b\\c&d\end{smallmatrix}\right)$,
forbidding $H_4$ imposes the single constraint
\[
 (a,b,c,d)\ne(1,0,1,1).
\]
The construction uses this exact negative constraint in four
orientations.  Two orientations push a value from level $i-1$ to
a cross-level cell, and two push it from that cell to the diagonal
at level $i$.  The plus chain propagates the value one; the minus
chain propagates zero.  Their terminal positions coincide.  The
gadget constraints therefore become inconsistent unless an entry
required by a selected frame has been changed.  The XOR choice of
physical representatives makes those required entries disjoint
across $m^2$ choices.
